%!TeX root=../main.tex

\chapter{Conclusão}
\label{cap:conclusao}
Os resultados do experimento, apesar de inconclusivos devido ao número reduzido de participantes, sugerem que a abordagem proposta apresentou um impacto positivo
no aprendizado dos fundamentos básicos do jogo. As métricas, em especial a variação na distribuição de probabilidades da seleção de exercícios de cada módulo,
evidenciam um progresso consistente por parte dos usuários segundo os critérios de desempenho estabelecidos pelo sistema. O questionário, por sua vez, ao coletar 
as impressões individuais dos participantes, destaca tanto os pontos nos quais o sistema contribuiu para o aprendizado quanto aqueles em que ainda precisa melhorar. 

O Shogi-sensei, portanto, cumpre com o objetivo de ensinar os fundamentos básicos de shogi para iniciantes ao adaptar as questões apresentadas com base no desempenho
do usuário e fornecer explicações para seus erros. Por outro lado, ainda se encontra em um estágio inicial de desenvolvimento, com o potencial de melhorar a 
efetividade do método, assim como a experiência do jogador.

Como trabalhos futuros, planeja-se focar principalmente em implementar as sugestões destacadas na Seção~\ref{sec:improvements}, como aumentar a consistência da
qualidade das explicações da inteligência artificial, melhorar a distribuição de peças nas questões e adicionar explicações básicas antes dos módulos sobre os
conceitos que serão abordados.


